\section{Anmerkungen}

\subsection{Landeslisten}

Das aktuelle Wahlprozedere sieht vor, dass die auf eine Partei entfallenden Bundestagsmandate
noch auf Landeslisten aufgeteilt werden. Landeslisten würde es nach dem hier vorgeschlagenen
Verfahren nicht mehr geben. Länder sollten bei der Bundestagswahl keine Rolle spielen. Die
regionale Vertretung findet bereits über die vorgeschlagene Verteilung der Wahlkreise statt.
Jede sich zur Wahl stellende Partei muss im ganzen Wahlgebiet, d.\,h.\ bei der Bundestagswahl im
ganzen Bundesgebiet, antreten.

Die Parteien dürfen intern Landeslisten führen, wenn sie das wünschen, um intern einen
Länderausgleich bei der Besetzung der Bundestagsliste durchzuführen. Das ist aber ihre
Entscheidung. Sie können über die nicht vom Bürger gewählte Hälfte ihrer Parlamentsmandate frei
verfügen.

\subsection{Fünf-Prozent-Hürde}

Das vorgestellte Verfahren enthält durch die Forderung, in allen Wahlkreisen mit drei Kandidaten
anzutreten, schon eine Hürde, die für kleine Parteien ein Problem sein könnte. Auch die drei
geforderten Spitzenkandidaten könnten für kleinere Parteien eine Hürde sein. Hinzu kommt, dass
sich die \enquote{natürliche Sperrklausel} (nur ganzzahlige Abgeordnete) durch das
\enquote{Zusammenstauchen} der Opposition noch stärker auswirkt.

Aus diesen Gründen denke ich, dass man auf eine explizite Fünf-Prozent-Hürde ggf.\ verzichten
könnte. In der Simulation wird sie nicht berücksichtigt. Falls sie aber doch politisch gewünscht
sein sollte, ließe sie sich leicht in das Verfahren einbauen.

\subsection{Kein CDU/CSU-Bündnis mehr}

Es ist aus staatstheoretischer Sicht unbedingt notwendig, dass alle sich zur Wahl stellenden
Parteien und Spitzenkandidaten im ganzen Wahlgebiet antreten. Jedes Parlaments- und
Regierungsmitglied ist dem gesamten Volk verpflichtet und darf nicht nur die Interessen einer
bestimmten Gruppe oder eines bestimmten Teilwahlgebiets vertreten. Zur Zeit hängt der Wahlerfolg
der CSU ausschließlich davon ab, wie sie beim bayerischen Wähler ankommt und nicht davon, was der
Rest Deutschlands von ihr hält. Dadurch haben wir durch die Hintertür doch wieder eine
Ländervertretung im Bundestag und die üblichen Schwierigkeiten zwischen Koalitionspartnern,
selbst wenn die CDU/CSU die absolute Mehrheit gewinnen würde. Regieren beinhaltet, unparteiische
Abwägungen zwischen verschiedenen Interessen zu treffen. Interessenvertretung ist Lobbyismus.
Interessenvertretung ist sehr notwendig und nichts Schlechtes in der Sache. Die Regierenden
können Interessen bei ihren Abwägungen nur berücksichtigen, wenn sie sie auch kennen. Die
Interessen müssen aber vor den Unparteiischen vertreten werden und nicht durch die Unparteiischen.
Eine Interessenvertretung der Länder geschieht bereits im Bundesrat. Dort ist auch Bayern
vertreten. Das muss reichen. Warum sollte Bayern privilegierter sein als alle anderen
Bundesländer?

Es ist eigentlich unverständlich und nur historisch zu erklären, dass das Bundesverfassungsgericht
zugelassen hat, dass eine zum Bundestag kandidierende Partei nicht in allen Ländern antritt.
So verkehrt, wie die CSU-Sonderlocke bereits jetzt ist, noch verkehrter würde sie nach dem hier
vorgeschlagenen Wahlverfahren. Das Wahlverfahren stellt ja extra sicher, dass Koalitionen nach
wie vor möglich, aber nicht mehr notwendig sind, um regieren zu können. Durch diese CDU/CSU-Kombination haben wir die Koalition aber garantiert, selbst bei nur zugesprochener absoluter
Mehrheit. Damit wäre das ganze Verfahren ad absurdum geführt. Wenn sowohl CDU als auch CSU in
allen Bundesländern antreten, können sie ruhig hinterher koalieren, wenn sie das wollen. Wenn die
CSU auf Bundesebene aber nicht in der Bedeutungslosigkeit verschwinden will, muss sie dazu
Konzepte entwickeln, die von ganz Deutschland und nicht nur von Bayern gut gefunden werden.
Vielleicht schafft sie das ja, denn offenbar laufen ja durchaus einige Dinge in Bayern besser als
im Rest der Republik. Der Länderfinanzausgleich spricht eine deutliche Sprache. Der Weg, sie auf
ganz Deutschland auszurollen, ist, selbst für ganz Deutschland Verantwortung zu übernehmen, indem
man in ganz Deutschland zur Wahl antritt. Vielleicht klappt es ja dann auch mit dem bayerischen
Bundeskanzler!

\subsection{Mehrere Christliche Parteien}

CDU/CSU könnten sich auch zu einer Partei zusammenschließen, um die Forderung nach bundesweit
antretenden Parteien zu erfüllen. Das wäre staatstheoretisch völlig in Ordnung, und vielleicht
werden sie das ja auch tun.

Allerdings wäre es wünschenswert, wenn auch nicht unbedingt notwendig, dass mehrere große
christliche Parteien bundesweit zur Wahl antreten, damit überzeugte christliche Wähler auch eine
echte politische Wahl bekommen und nicht nur nach Religionszugehörigkeit wählen müssen. Zurzeit
wählen viele überzeugte Christen CDU/CSU aufgrund des \enquote{C}s im Namen, obwohl sie mit
deren politischen Ansichten nicht unbedingt übereinstimmen. Wenn mehrere christliche Parteien mit
unterschiedlichen politischen Ansichten zur Auswahl stünden, wäre dies sicherlich förderlich für
die Demokratie.
